\section{Adaptive scales with one Gaussian side projection}
\label{l:side:section}

We now use a different information balance.  A finite message carries
little information when the prior has sufficiently small mass near
lower-dimensional affine planes.  For a general bounded-density prior,
a maximizing radius splits it into cells with that property.  An independent
Gaussian projection lets us bound the conditional entropy of the cell
label.  We charge this revelation cost in each block's information
estimate, keeping the same projection throughout the transcript argument.
It is an analysis
variable; the learner still receives only its prescribed exact samples.

All logarithms and entropies in this section use natural units.  Let
$\mu$ be uniform probability on $S^{d-1}$.  For sufficiently large $d$,
fix
\begin{equation}\label{l:side:parameters}
 k_*=p=\lfloor d/32\rfloor,\qquad q=p-1,\qquad
 a=d/8,\qquad \alpha=d/4,\qquad \ell=\lfloor d/3\rfloor.
\end{equation}
Every constant denoted by $C$ is positive and absolute.

We will use Lemma~\ref{l:caps-endpoint} also in lower ambient
dimensions.  If $\mu_n$ is uniform probability on $S^{n-1}$, $n\ge2$,
that lemma gives, for $s\in S^{n-1}$,
\begin{equation}\label{l:side:caps}
 \mu_n(B(s,r))\le r^{n-1}\quad(0<r\le1/2),\qquad
 \mu_n(B(s,r))\le(2r)^{n-1}\quad(r>0).
\end{equation}

\begin{lemma}[A finite message from a locally diffuse prior]
\label{l:side:local-information}
Let $\nu\ll\mu$ be a probability measure supported in $B(c,R)$, where
$c\in\mathbb R^d$ and $R>0$.  Suppose that for $\nu$-almost every $s$, simultaneously for all
$r>0$,
\begin{equation}\label{l:side:local-mass}
 \nu(B(s,r))\le K(r/R)^\alpha,\qquad K\ge1.
\end{equation}
Draw $S\sim\nu$ and an independent standard Gaussian $k$-by-$d$ matrix
$A$, where $1\le k\le k_*$.  A measurable probability kernel of
$(A,AS)$ produces a message $Z$ in an alphabet of cardinality $w\ge1$.
Then
\begin{equation}\label{l:side:local-bound}
 I(S;Z)\le C(d+\log K)+\frac{\log w}{p}.
\end{equation}
\end{lemma}

\begin{proof}
We first control the volume of the simplex spanned by independent
points from $\nu$.  The use of Gaussian replicas and Gram--Schmidt
factors is related to the projection argument of Sharan, Sidford, and
Valiant~\cite[Section~7, Lemma~12]{SSV2019}.  The geometric moment and
the exact-message information estimate needed here are proved below.

Let $V$ be an affine subspace of dimension $m\le p-2$, and let
$0<u\le1$.  Points of $B(c,R)$ within distance $uR$ of $V$ project into
$V\cap B(c,2R)$.  A maximal separated set covers this set by at most
$(C/u)^m$ balls in $V$ of radius $uR$; when $m=0$, one ball suffices.
The corresponding points in the original tube lie in ambient balls of
radius $2uR$.  Such a ball of positive $\nu$-mass contains a point where
\eqref{l:side:local-mass} holds, and it is contained, on the support of
$\nu$, in a radius-$4uR$ ball about that point.  Hence
\begin{equation}\label{l:side:tube}
 \nu\{s:\operatorname{dist}(s,V)\le uR\}
 \le KC^d u^{\alpha-m}.
\end{equation}
The parameters give $\alpha-m-a\ge3d/32$.  Integrating the tail of
$(R/\operatorname{dist}(S,V))^a$, and using
\eqref{l:side:tube} for values of this variable above one, gives
\begin{equation}\label{l:side:height-moment}
 \mathbb E_\nu (R/\operatorname{dist}(S,V))^a\le KC^d.
\end{equation}
Indeed the tail integral is bounded by
$1+KC^d a/(\alpha-m-a)$.  The tube bound also proves that $V$ has
zero $\nu$-mass.

Draw $(S_1,\ldots,S_p)\sim\nu^{\otimes p}$.  Let $F$ be the matrix
with columns $S_i-S_1$, $2\le i\le p$, and set
$v=\sqrt{\det(F^{\mathsf T}F)}$.  Gram--Schmidt writes $v$ as the
product of the distances from $S_i$ to
$\operatorname{aff}(S_1,\ldots,S_{i-1})$.  The preceding zero-mass
statement shows successively that these distances are positive almost
surely.  Integrating the last point first and repeating
\eqref{l:side:height-moment} proves
\begin{equation}\label{l:side:gram-moment}
 \mathbb E_{\nu^{\otimes p}}(R^q/v)^a\le(KC^d)^q.
\end{equation}

We use this moment in a normalized comparison of posterior replicas.
Temporarily add an independent $N(0,\sigma^2I_k)$ vector $\eta$ to the
labels, where $0<\sigma\le R$, and set $b=AS+\eta$.  Apply the given
message kernel to $(A,b)$ and call its output $Z_\sigma$.  Write $P_A$
for the Gaussian matrix law.  With $\varphi_\sigma$ denoting the noise density, write
\[
 f_\sigma(b\mid A)=\int\varphi_\sigma(b-As)\,d\nu(s).
\]
Sample $(A,b,Z_\sigma)$ from this experiment.  Conditional on these
variables, sample $p$ independent points from the posterior of $S$
given $(A,b)$; their conditional law does not depend further on
$Z_\sigma$.  Let $P_\sigma$ be the resulting marginal on
$(S_1,\ldots,S_p,A,b)$ and $Q_\sigma$ its point-tuple marginal.
Each pair $(S_i,Z_\sigma)$ has the experiment's original marginal.
Put $D_Q=D(Q_\sigma\|\nu^{\otimes p})$, initially allowing the value
$+\infty$.  Once $D_Q$ is finite, relative entropy to a product bounds
the sum of the marginal relative entropies: their difference is the
relative entropy of the joint law to the product of its own marginals.
Apply this after conditioning on the common message and use the matching
single-point marginals.  The conditional chain rule then gives
\begin{align}
 pI(S;Z_\sigma)
 &\le\mathbb E_{Z_\sigma}
   D(Q_{\sigma,S_1,\ldots,S_p\mid Z_\sigma}\|\nu^{\otimes p})\notag\\
 &=D_Q+I((S_1,\ldots,S_p);Z_\sigma)
 \le D_Q+\log w.
 \label{l:side:replica-message}
\end{align}
It therefore suffices to prove $D_Q/p\le C(d+\log K)$, uniformly in
$\sigma$.  We will prove finiteness and this bound by comparing the
posterior tuple law with a normalized reference whose tuple marginal is
exactly $\nu^{\otimes p}$.

Relative to $\nu^{\otimes p}\otimes P_A\otimes db$, the density of
$P_\sigma$ is
\begin{equation}\label{l:side:auxiliary-density}
 \frac{\prod_{i=1}^p\varphi_\sigma(b-AS_i)}
      {f_\sigma(b\mid A)^q}.
\end{equation}

For a fixed tuple, the integral of the numerator over $(A,b)$ is
\begin{equation}\label{l:side:normalizer}
 c_\sigma(F)=(2\pi)^{-kq/2}
 \det\!\left(F^{\mathsf T}F+
       \sigma^2(I_q+\mathbf1\mathbf1^{\mathsf T})\right)^{-k/2}.
\end{equation}
Here $\mathbf1\in\mathbb R^q$ is the all-ones column.
To verify this formula, take $p$ independent noisy labels with the
same matrix $A$.  Integrating their product density over their common
value $b$ evaluates at zero the density of the $q$ differences from
the first label.  In one row, the projected signal differences have
covariance $F^{\mathsf T}F$, and the noise differences have covariance
$\sigma^2(I_q+\mathbf1\mathbf1^{\mathsf T})$.  The rows are independent,
which proves \eqref{l:side:normalizer}.  On full-rank tuples,
\begin{equation}\label{l:side:normalizer-volume}
 c_\sigma(F)\le(2\pi)^{-kq/2}v^{-k}.
\end{equation}

Define a reference probability law $R_\sigma$ by the density
\begin{equation}\label{l:side:reference-density}
 \frac{\prod_{i=1}^p\varphi_\sigma(b-AS_i)}{c_\sigma(F)}
 \quad\text{relative to }\nu^{\otimes p}\otimes P_A\otimes db.
\end{equation}
The denominator is exactly the integral for each fixed tuple.  Thus
the tuple marginal of this reference law is $\nu^{\otimes p}$, rather
than an unnormalized multiple of that law.

All logarithms in this comparison are integrable at fixed $\sigma>0$.
The tuple is supported in $B(c,R)^p$, and
\[
 \sigma^2I_q\preceq F^{\mathsf T}F+
       \sigma^2(I_q+\mathbf1\mathbf1^{\mathsf T})
 \preceq(4qR^2+p\sigma^2)I_q.
\]
Hence $\log c_\sigma$ is bounded as the tuple varies.  Also
$\log f_\sigma$ is at most $-\frac{k}{2}\log(2\pi\sigma^2)$ and at
least this value minus
\[
 \frac{(\|b-Ac\|+\|A\|_{\mathrm{op}}R)^2}{2\sigma^2}.
\]
The latter term is integrable under the noisy experiment.  Data processing
therefore proves $D_Q<\infty$ and gives
\begin{equation}\label{l:side:kl-comparison}
 D_Q\le D(P_\sigma\|R_\sigma)
 =\mathbb E_{P_\sigma}\log c_\sigma(F)
       -q\mathbb E_{P_\sigma}\log f_\sigma(b\mid A).
\end{equation}
To bound the second term, compare $f_\sigma(\cdot\mid A)$ with the
$N(Ac,2R^2I_k)$ density.  Since
$\mathbb E\|b-Ac\|^2\le kR^2+k\sigma^2\le2kR^2$, nonnegative
relative entropy yields
\[
 -\mathbb E\log f_\sigma(b\mid A)
 \le k\log R+\frac{k}{2}\log(4\pi e).
\]
Together with \eqref{l:side:normalizer-volume}, this gives, initially
allowing the expectation on the right to be extended,
\begin{equation}\label{l:side:kl-volume}
 D_Q\le k\mathbb E_{Q_\sigma}\log(R^q/v)+Ckq.
\end{equation}

Here the needed expectation is in fact finite.  Put
$X=a\log(R^q/v)$.  Hadamard's inequality gives $v\le(2R)^q$, so
$X\ge-aq\log2$.  For bounded $X'$ the probability law proportional to
$e^{X'}\nu^{\otimes p}$ and nonnegative relative entropy give
\[
 \mathbb E_{Q_\sigma}X'
 \le D_Q+\log\mathbb E_{\nu^{\otimes p}}e^{X'}.
\]
Apply this to $X'=\min(X,n)$ and let $n\to\infty$.  Monotone convergence
after adding the fixed lower bound, together with
\eqref{l:side:gram-moment}, proves both finiteness and
\[
 a\mathbb E_{Q_\sigma}\log(R^q/v)
 \le D_Q+q(\log K+Cd).
\]
Because $k/a\le1/4$, substitution into \eqref{l:side:kl-volume} and
absorption of $(k/a)D_Q$ give
\begin{equation}\label{l:side:kl-absorption}
 D_Q/p\le C(d+\log K).
\end{equation}

Combining \eqref{l:side:kl-absorption} with
\eqref{l:side:replica-message} proves \eqref{l:side:local-bound} for the
noisy experiment, uniformly over $0<\sigma\le R$.

It remains to pass this inequality to the given exact-label rule.  This
step uses the rule itself, rather than a comparison of the statistical
difficulties of noisy and noiseless experiments.  For every rank-$k$
matrix with $k<d$, the image of $\mu$ has a Lebesgue density.  One direct
proof uses a normalized standard Gaussian: the squared length of the
first $k$ coordinates is a beta variable with parameters
$k/2,(d-k)/2$, obtained from the ratio of independent chi-square
variables, and its direction is uniform.  This gives a density on the
unit $k$-ball.  Orthogonal changes of coordinates and an invertible
range map give the claim for any rank-$k$ matrix.  Absolute continuity
$\nu\ll\mu$ preserves it.  Since a Gaussian matrix has full row rank
almost surely and itself has a Lebesgue density, $(A,AS)$ has a joint
Lebesgue density.

Let $g_z(A,b)$ be the output probability of message $z$.  For any bounded
measurable function $g$ and a random vector $X$ with a Lebesgue density,
\begin{equation}\label{l:side:weighted-translation}
 \mathbb E|g(X+h)-g(X)|\longrightarrow0\quad(h\to0).
\end{equation}
To see this, restrict the input density to a bounded cube and to values
below a fixed bound, losing arbitrarily little probability.  For small
$h$, cut $g$ off on a larger cube.  The remaining integral is bounded
by a constant times the Lebesgue $L^1$ distance between this cutoff and
its translate.  That distance tends to zero, since continuous compactly
supported functions are dense in $L^1$ and their translates converge
uniformly on a common compact set.  This proves
\eqref{l:side:weighted-translation}.  Apply it to shifts in the label
coordinates, then to $h=(0,\sigma G)$ with $G$ standard Gaussian by
bounded convergence.  The finite alphabet gives
\begin{equation}\label{l:side:kernel-limit}
 \sum_z\mathbb E|g_z(A,AS+\eta)-g_z(A,AS)|\longrightarrow0.
\end{equation}
Couple the two message probability vectors at each $(S,A,\eta)$ by their
common masses, so that their disagreement probability $\delta$ is half
the expected sum in \eqref{l:side:kernel-limit}.  This is a measurable
coupling because the alphabet is finite.  Write
$h_2(t)=-t\log t-(1-t)\log(1-t)$, with $0\log0=0$, for binary entropy.
The disagreement indicator bounds each of the marginal and conditional entropy differences by
$h_2(\delta)+\delta\log w$: on disagreement there are at most $w$
possible values, and binary entropy is concave when conditional
disagreement probabilities are averaged.  Consequently
\[
 |I(S;Z_\sigma)-I(S;Z)|
 \le2\{h_2(\delta)+\delta\log w\}\longrightarrow0.
\]
For $w=1$ both messages are constant and the assertion is immediate.
The uniform noisy bound therefore proves \eqref{l:side:local-bound}
for the exact labels $AS$.
\end{proof}

\begin{lemma}[Adaptive cells predicted by independent side data]
\label{l:side:adaptive}
Let $\nu$ be a probability measure with $d\nu/d\mu\le B<\infty$.
Draw $S\sim\nu$, an independent
standard Gaussian $k$-by-$d$ matrix $A$ with $1\le k\le k_*$, and a
message $Z$ in an alphabet of size $w\ge1$ from a measurable kernel of
$(A,AS)$.  Independently of the signal, the matrix, and the message
randomness, draw a standard Gaussian $\ell$-by-$d$ matrix $W$, and put
$D_*=(W,WS)$.  Then
\begin{equation}\label{l:side:adaptive-bound}
 I(S;Z\mid D_*)\le
 C\left[d+\log\left(C\left(2+\frac{\log B}{d}\right)\right)\right]
 +\frac{\log w}{p}.
\end{equation}
\end{lemma}

\begin{proof}
Since $\nu$ is a probability, $B\ge1$.  Put
$m_s(r)=\nu(B(s,r))$ and $R_j=2^{1-j}$ for integers $j\ge0$.
Choose $J(s)$ to be the least index maximizing
$2^{\alpha j}m_s(R_j)$, and write $J=J(S)$.  The score at $j=0$ is one because $R_0=2$
covers the sphere.  The second cap bound in \eqref{l:side:caps} gives
\[
 2^{\alpha j}m_s(R_j)
 \le B\,2^{2(d-1)-j(d-1-\alpha)}\longrightarrow0.
\]
The ball masses are measurable in $s$, by integrating the Borel indicator
of $\|s-u\|\le R_j$.  Thus the least maximizer is measurable and
\begin{equation}\label{l:side:scale-entropy}
 0\le J(s)\le\frac{\log_2 B+2(d-1)}{d-1-\alpha},\qquad
 H(J)\le\log\left(C\left(2+\frac{\log B}{d}\right)\right).
\end{equation}
The entropy bound follows from the displayed finite range.  For
$R=R_{J(s)}$, compare the maximum with the least dyadic radius at least
$r$ when $r\le2$, and with $R_0=2$ when $r>2$.  This proves
\begin{equation}\label{l:side:adaptive-growth}
 m_s(r)\le2^\alpha m_s(R)(r/R)^\alpha\qquad(r>0).
\end{equation}

At each scale choose a finite maximal $R_j$-separated set of centers on
the sphere.  Its $R_j$-balls cover the sphere.  Enumerate the centers and
assign a point to the first center within distance $R_j$; this is a
measurable rule.  Let $U=(J,i)$ record the scale and assigned center.
Its range is finite by \eqref{l:side:scale-entropy} and the finiteness of each net.
For each cell $u=(j,i)$ of positive probability, write $c_u$ for its
center, $b_u=\mathbb P(U=u)$, and
$m_u^+=\nu(B(c_u,2R_j))$.

At one scale, a point lies in at most $5^d$ of the enlarged balls
$B(c_u,2R_j)$.  Indeed the disjoint interiors of the radius-$R_j/2$
balls about those centers lie in a ball of radius $5R_j/2$ about the
point.  Volume comparison proves the count, and integration gives
$\sum_{u=(j,i)}m_u^+\le5^d$.  Write $p_j=\mathbb P(J=j)$.  Concavity of
the logarithm with weights $b_u/p_j$ at each positive-mass scale yields
\begin{equation}\label{l:side:overlap-entropy}
 \sum_u b_u\log\frac{m_u^+}{b_u}
 \le\sum_{j:p_j>0}p_j\log\frac{5^d}{p_j}
 =d\log5+H(J).
\end{equation}

The conditional prior $\nu_u$ is supported in $B(c_u,R_j)$ and remains
absolutely continuous with respect to $\mu$.  If $s$ is in that cell,
then $m_s(R_j)\le m_u^+$.  Dividing \eqref{l:side:adaptive-growth} by
$b_u$ gives, for $\nu_u$-almost every $s$ and all $r>0$,
\[
 \nu_u(B(s,r))\le
 \frac{2^\alpha m_u^+}{b_u}(r/R_j)^\alpha.
\]
Here $m_u^+\ge b_u$, so the factor is at least one.  The matrix $A$
remains independent of $(S,U)$, and the same message kernel is used.
Apply Lemma~\ref{l:side:local-information} for each $u$, average, and use
\eqref{l:side:overlap-entropy}.  We obtain
\begin{equation}\label{l:side:information-in-cell}
 I(S;Z\mid U)\le C(d+H(J))+\frac{\log w}{p}.
\end{equation}

We have bounded the information inside each cell.  We now pay for its
label in the presence of the side data.  For a known scale $j$, predict
the center index from $(W,y)$ by the normalized probabilities
\begin{equation}\label{l:side:predictor}
 \pi_{j,i}(W,y)=
 \frac{b_{(j,i)}\exp(-\|Wc_{(j,i)}-y\|^2/(2R_j^2))}
 {\displaystyle\sum_{i'}b_{(j,i')}
       \exp(-\|Wc_{(j,i')}-y\|^2/(2R_j^2))}.
\end{equation}
Discard zero-probability cells.  At an occurring scale the denominator
is positive and the set of indices is finite.  Nonnegative relative
entropy between the conditional center distribution and these
probabilities bounds $H(i\mid D_*,J)$ by the expected negative logarithm
of the true cell's prediction.

Fix a true point $s$ with $J(s)=j$ and put $R=R_j$.  For its center,
\[
 \mathbb E_W\frac{\|W(c_U-s)\|^2}{2R^2}\le\ell/2.
\]
For any candidate center $c$, direct Gaussian integration gives
\[
 \mathbb E_W\exp(-\|W(c-s)\|^2/(2R^2))
 =\left(1+\frac{\|c-s\|^2}{R^2}\right)^{-\ell/2}.
\]
Centers within distance $2R$ of $s$ have total cell probability at most
$m_s(3R)$, since their cells are within an additional distance $R$.
Centers at distances in $(2^tR,2^{t+1}R]$, $t\ge1$, have total cell
probability at most $m_s(2^{t+2}R)$ and expected kernel at most
$2^{-t\ell}$ per unit cell probability.  The expected denominator in
\eqref{l:side:predictor}, evaluated at $y=Ws$, is therefore at most
\begin{align}
 m_s(3R)+\sum_{t\ge1}2^{-t\ell}m_s(2^{t+2}R)
 &\le m_s(R)\left(6^\alpha+
       2^{3\alpha}\sum_{t\ge1}2^{-t(\ell-\alpha)}\right)\notag\\
 &\le e^{Cd}m_s(R).
 \label{l:side:shells}
\end{align}
The first inequality uses \eqref{l:side:adaptive-growth}; the series
converges because $\ell-\alpha\ge1$ for large $d$.

Jensen's inequality for the logarithm of this denominator, together
with the true-center squared term, bounds the expected prediction loss
at $s$ by $Cd+\log(m_s(R)/b_U)$.  The logarithms are integrable: the
denominator is at least the true term, and its negative logarithm is
bounded by $-\log b_U$ plus a Gaussian squared term.  Averaging and using
$m_S(R_J)\le m_U^+$ in \eqref{l:side:overlap-entropy} gives
\begin{equation}\label{l:side:cell-given-side}
 H(U\mid D_*)
 \le H(J)+Cd+\mathbb E\log\frac{m_S(R_J)}{b_U}
 \le Cd+2H(J).
\end{equation}

Finally $U$ is a function of $S$, and $Z$ is conditionally independent
of $D_*$ given $(S,U)$.  Since $Z$ is finite,
\begin{align*}
 I(S;Z\mid D_*,U)
 &=H(Z\mid D_*,U)-H(Z\mid S,U)\\
 &\le H(Z\mid U)-H(Z\mid S,U)=I(S;Z\mid U).
\end{align*}
The chain rule then gives
$I(S;Z\mid D_*)\le H(U\mid D_*)+I(S;Z\mid U)$.
Combining \eqref{l:side:scale-entropy},
\eqref{l:side:information-in-cell}, and
\eqref{l:side:cell-given-side} proves \eqref{l:side:adaptive-bound}.
The cell has been revealed only inside this calculation, with its full
conditional entropy charged; no local regularity of the original
state posterior was assumed.
\end{proof}

\subsection{The full state transcript and the precision endpoint}
\label{l:side:endpoint}

We now give a second proof of Corollary~\ref{l:grid:precision}.
The information bound will be conditional on the same side data
$D_*$ at every block, and its endpoint will account for the residual
sphere left by those data.

\begin{proof}[Proof of Corollary~\ref{l:grid:precision} using side data]
Take $S\sim\mu$ and use the reductions following
Lemma~\ref{l:prior-ac}.  First fix the shared rule randomness with
uniform-prior success greater than $3/5$, apply Lemma~\ref{l:borel},
and then fix the independent initial state with success at least
$3/5$.  Fresh transition and terminal output randomness may remain:
the two lemmas permit finite message kernels, and the calculation below
integrates the terminal output law.

Put $L=\log(1/\epsilon)$ with $0<\epsilon\le1/10$.
If $T=0$, the fixed initial state and terminal rule give an output
independent of $S$, whose success is at most $\epsilon^{d-1}<3/5$ by
\eqref{l:side:caps}.  Thus the successful runs considered below have
$T\ge1$.

It suffices to consider $T<dL$.  There are at most $(T+1)2^M$ terminal pairs consisting
of the stopping index and state.  For a fixed pair $h$, the joint
subprobability law $\lambda_h$ of $S$ and the event of reaching $h$
satisfies $\lambda_h\le\mu$.  The unit output conditional on $h$ uses
only its terminal law and fresh randomness, so is independent of $S$.
Angular success implies Euclidean error at most $\epsilon$, and
\eqref{l:side:caps} applies since $\epsilon\le1/10$.  Integrating the
terminal output law and summing over pairs gives
\begin{equation}\label{l:side:output-count}
 \frac35\le(T+1)2^M\epsilon^{d-1},\qquad
 (d-1)L\le M\log2+\log(T+1)+O(1).
\end{equation}
Using $T<dL$ and $\log(dL+1)\le\log(d+1)+L$ yields
$(d-2)L\le M\log2+O(\log d)$.  Along a sequence $M=o(d^2)$ this
implies $L=o(d)$ in the present short-run case and eventually $T\le d^2$.

After an early stop, preserve its state together with its terminal index
and ignore the remaining fresh samples.  The active states and all
terminal tags together give the global width bound
\begin{equation}\label{l:side:padded-width}
 w=(T+2)2^M,\qquad \log w=M\log2+\log(T+2)=o(d^2).
\end{equation}
The output law at a terminal tag is the original terminal law.  Split
the $T$ padded samples into $N=\lceil T/k_*\rceil$ consecutive blocks
of length at most $k_*$.  Let $\Pi_i$ record every state at the first
$i$ block ends, with $\Pi_0$ empty.  This is the full state transcript,
not just its final state.

Independently draw one matrix $W$ with $\ell$ Gaussian rows and retain
the same $D_*=(W,WS)$ for all blocks.  For every positive-mass history
$h$ of $\Pi_{i-1}$, its conditional signal law satisfies
\begin{equation}\label{l:side:history-density}
 \frac{d\nu_h}{d\mu}(s)
 =\frac{\mathbb P(\Pi_{i-1}=h\mid S=s)}
        {\mathbb P(\Pi_{i-1}=h)}
 \le B_h:=\frac1{\mathbb P(\Pi_{i-1}=h)}.
\end{equation}
Given $h$, the next block matrix is still fresh Gaussian and independent
of $S$, and $W$ remains independent of both.  The incoming state is known
from $h$.  The persistent-memory restriction makes the successor state
a kernel of that state and the block's complete exact data, with alphabet
size at most $w$.  It has no access to $D_*$.  Lemma~\ref{l:side:adaptive}
therefore applies conditionally on each history.

The posterior-density cost is small after averaging.  Since $\Pi_{i-1}$
is a tuple of at most $N$ states from an alphabet of size $w$,
\[
 \mathbb E_h\log B_h=H(\Pi_{i-1})\le N\log w.
\]
Jensen's inequality gives
\begin{align*}
 \mathbb E_h\log\left(C\left(2+\frac{\log B_h}{d}\right)\right)
 &\le\log\left(C\left(2+\frac{N\log w}{d}\right)\right)
 =O(\log d).
\end{align*}
Here $N=O(d)$ follows from $T\le d^2$ and $k_*\ge d/64$ for large $d$;
also $(\log w)/p=o(d)$ by \eqref{l:side:padded-width}.  Summing the
conditional information bounds for successive states by the chain rule
gives an absolute $C_1$ such that
\begin{equation}\label{l:side:transcript-upper}
 I(S;\Pi_N\mid D_*)\le C_1dN.
\end{equation}
The local cell $U$ used in each application is not an uncharged part of
$\Pi_N$: its revelation was paid by \eqref{l:side:cell-given-side} inside
that block's conditional bound.

For a lower bound on the same information, condition on $D_*$.  Almost
surely $W$ has rank $\ell$.  Its labels specify $P_WS$, where $P_W$ is
the orthogonal projector onto its row space.  The remaining component
has uniform direction in the nullspace, of dimension $d-\ell$, and radius
\[
 \rho=\sqrt{1-\|P_WS\|^2}.
\]
This conditional law follows by writing $S$ as a normalized standard
Gaussian: the nullspace direction is independent of its length and of
the row-space component.  The event $\mathcal G=\{\rho\ge1/2\}$ is
determined by $D_*$, and Markov's inequality gives
\begin{equation}\label{l:side:good-radius}
 \mathbb P(\mathcal G^c)
 \le\frac43\mathbb E\|P_WS\|^2
 =\frac{4\ell}{3d}\le\frac49.
\end{equation}
On $\mathcal G$, the mass of the conditional sphere within Euclidean
distance $\epsilon$ of any fixed estimate is at most
$(4\epsilon)^{d-\ell-1}$.  If the ball meets that sphere, choose a point
of intersection.  The rest of the intersection lies within $2\epsilon$
of it.  After rescaling the residual sphere by $\rho\ge1/2$, the radius
is at most $4\epsilon\le2/5<1/2$, so the first bound in
\eqref{l:side:caps} applies in ambient dimension $d-\ell$.

The final output $\widehat S$ is conditionally independent of $S$ given
$(\Pi_N,D_*)$, because the transcript contains its terminal tag.
For each value of $D_*$, compare the actual joint law of
$(S,\Pi_N,\widehat S)$ with the product of its conditional $S$ marginal
and conditional $(\Pi_N,\widehat S)$ marginal.  On $\mathcal G$ the
success probability under this product is at most
$(4\epsilon)^{d-\ell-1}$.  Binary relative entropy is at least the true
event probability times the negative logarithm of its product-law
probability, minus $\log2$.  Integrating over $\mathcal G$, and using
nonnegativity elsewhere, gives
\begin{align}
 I(S;\Pi_N\mid D_*)
 &=I(S;\Pi_N,\widehat S\mid D_*)\notag\\
 &\ge\left(\frac35-\frac49\right)(d-\ell-1)
       \log\frac1{4\epsilon}-\log2
 \ge c_0dL
 \label{l:side:transcript-lower}
\end{align}
for an absolute $c_0>0$ and sufficiently large $d$.  The coefficient
$3/5-4/9$ accounts for the success mass outside the failed-radius event.
For the last inequality, $d-\ell-1$ is a positive multiple of $d$ and
$L-\log4\ge(1-\log4/\log10)L$.

Combining \eqref{l:side:transcript-upper} and
\eqref{l:side:transcript-lower}, and using
$N=\lceil T/k_*\rceil\le T/k_*+1$, gives the exact intermediate bound
\begin{equation}\label{l:side:precision-time}
 T\ge k_*\left(\frac{c_0}{C_1}L-1\right).
\end{equation}
We retain a separate linear endpoint to absorb its additive constant.
If $T<d/3$, reveal all $T$ pre-generated rows and exact labels, even after
an early stop.  The same normalized-Gaussian argument gives a uniform
conditional sphere of ambient dimension $d-T$.  Its radius is at least
$1/2$ except on an event of probability at most $4T/(3d)<4/9$.
The preceding conditional cap argument therefore bounds the success of
even a full-data estimator by
\[
 \frac49+(4\epsilon)^{d-T-1}<\frac35
\]
for sufficiently large $d$, a contradiction.  Thus $T\ge d/3$.
When $L\ge2C_1/c_0$, \eqref{l:side:precision-time} implies $T\ge cdL$;
for smaller $L$, the linear bound gives the same conclusion after reducing
the absolute constant.  The case $T\ge dL$ needed no reduction.
Every learner observation throughout this argument is the exact inner
product; the auxiliary noise appeared only in the proved passage to
the exact-message inequality.
\end{proof}
